\begin{frame}{Where this deck goes --- one concept per frame, then the evidence}
\vspace{-1mm}
\begin{columns}[T,onlytextwidth]
\begin{column}{0.52\textwidth}
\textbf{\color{acc}Decides WHEN} \tsub{— the schedule model}\\[1mm]
{\scriptsize
\begin{tabular}{@{}rp{0.78\linewidth}@{}}
\tsub{2} & \textbf{What a schedule IS} — seven fields, op-class, presence, rings\\
\tsub{3} & \textbf{Placement} — retiming a schedule into a pipeline\\
\tsub{4} & \textbf{Obligations} — the correctness argument\\
\end{tabular}}

\vspace{2.5mm}
\textbf{\color{acc}Resolves WHERE} \tsub{— the dataflow IR}\\[1mm]
{\scriptsize
\begin{tabular}{@{}rp{0.78\linewidth}@{}}
\tsub{5} & \textbf{Why a second IR} — and what a tree cannot answer\\
\tsub{6} & \textbf{The analyses} — generations, tokens, hazards, fences\\
\tsub{7} & \textbf{Passes and verification} — and the emission boundary\\
\end{tabular}}

\vspace{2.5mm}
\textbf{\color{acc}Decides HOW} \tsub{— the leaves}\\[1mm]
{\scriptsize
\begin{tabular}{@{}rp{0.78\linewidth}@{}}
\tsub{8} & \textbf{Address geometry} — the one thing the model may not know\\
\end{tabular}}

\vspace{2.5mm}
\textbf{\color{acc}Evidence}\\[1mm]
{\scriptsize
\begin{tabular}{@{}rp{0.78\linewidth}@{}}
\tsub{9} & \textbf{What shipped}, and how it is validated\\
\tsub{10} & \textbf{Known gaps}, and the bet this design makes\\
\end{tabular}}
\end{column}

\begin{column}{0.46\textwidth}
\vspace{1mm}
\begin{tikzpicture}
  \node[box](w){WHEN};
  \node[box,right=4mm of w](r){WHERE};
  \node[gbox,right=4mm of r](h){HOW};
  \draw[ar](w)--(r); \draw[ar](r)--(h);
\end{tikzpicture}\\
\tsub{Each layer is forbidden to know the next one's answer.}

\vspace{4mm}
\begin{tikzpicture}
\node[wbox,text width=0.95\columnwidth,align=left]{%
\scriptsize\textbf{This is the architecture deck.}\\[3pt]
\raggedright
One frame per concept, and no frame assumes a term the earlier ones have not defined.\\[4pt]
The \key{appendix} holds the deep dives — the derivations, the post-mortems and the
per-feature matrices — each attached to a spine claim. They are there to be \emph{opened when
challenged}, not presented in order.\\[4pt]
\tsub{Two in-tree implementation reports carry the full detail; this deck carries the shape.}};
\end{tikzpicture}
\end{column}
\end{columns}
\end{frame}
